\documentclass[11pt]{article}
\usepackage[utf8]{inputenc}
\usepackage[T1]{fontenc}
\usepackage{lmodern}
\usepackage[margin=2.4cm]{geometry}
\usepackage{xcolor}
\usepackage[hidelinks]{hyperref}
\newcommand{\resp}[1]{\par\noindent\textcolor{blue!60!black}{\textbf{Response:}} #1\par\medskip}
\setlength{\parindent}{0pt}
\setlength{\parskip}{0.4em}

\begin{document}
\begin{center}
{\large\textbf{Response to Reviewer --- Round 2}}\\[0.3em]
\textit{Environmental Monitoring and Assessment}\\[0.2em]
Manuscript: ``A Transferable Protocol for Satellite-Based Water Transparency
Monitoring in High-Altitude Oligotrophic Lakes: Harmonized Sentinel-2/Landsat
Retrieval with Rigorous Validation and Uncertainty Quantification in Lake Titicaca
(2013--2024)''
\end{center}

We sincerely thank the Reviewer for the very positive assessment and for the
recognition that the previous revisions substantially improved the manuscript's
alignment with the scope of \textit{Environmental Monitoring and Assessment}. We
have now addressed all five minor points; each change is incorporated in the
revised manuscript.

\hrulefill

\textbf{1. Abstract --- improve flow; the first sentence was too long.}
\resp{We rewrote the Abstract along the lines kindly suggested by the Reviewer. The
opening is now a short, direct sentence (``High-altitude Andean lakes are critical
freshwater reserves that remain poorly monitored by sparse field campaigns''), the
protocol framing leads, and the ``operational envelope'' and the ``provided local
recalibration is performed'' caveat are stated explicitly. We retained a single
compact sentence with the key quantitative results ($R^2$, RMSE, conformal
coverage), which we judged appropriate for a data-driven monitoring journal; the
Abstract is now tighter (173 words) and reads more fluently.}

\textbf{2. Figure 1 --- make it cleaner/more professional; legible in black and
white.}
\resp{Figure 1 has been replaced with a professionally redesigned, high-resolution
version (clean rounded boxes, consistent typography, clearer directional arrows).
The diagram follows a single top-to-bottom flow with fully labelled stages, so the
sequence remains unambiguous in grayscale/print, where the boxes are distinguished
by position and text rather than colour alone.}

\textbf{3. Section 5.1 (Operational implications) --- (a) note the real cloud
limitation on cadence; (b) make the local-recalibration recommendation more
emphatic.}
\resp{(a) We added an explicit, honest statement that persistent cloud cover ---
frequent during the austral wet season (December--March) --- reduces the number of
usable scenes, so the realized update frequency is lower than the nominal revisit
and is seasonally variable, and that this seasonal availability should be reported
to managers alongside the maps. (b) We added an emphatic closing sentence to the
section stating that the protocol \emph{must not} be transferred to a new lake or an
unsampled optical sub-basin without prior local recalibration, since applying it
blindly would risk confidently wrong maps precisely where managers most need
reliable information.}

\textbf{4. Language and flow --- split long sentences; polish.}
\resp{We split the long ``Framed as an environmental-monitoring tool\ldots''
sentence in the Introduction into two clearer sentences, and made several minor
edits to shorten long sentences and smooth paragraph transitions in the Discussion.}

\textbf{5. Conclusions --- make the management-utility message more direct.}
\resp{We strengthened the closing of the Conclusions with a direct statement that
the protocol gives Andean lake managers a low-cost, continuously updated instrument
to detect clarity loss early and act before degradation becomes difficult to
reverse, directly serving the environmental monitoring and assessment of these
transboundary freshwater reserves.}

\hrulefill

We are grateful for the Reviewer's constructive guidance across both rounds, which
has clearly strengthened the manuscript. We hope the revised version is now suitable
for publication in \textit{Environmental Monitoring and Assessment}.

\end{document}
